% ===========================================================================
\chapter{The agent loop}\label{ch:loop}

We have a conversation, a model and tools. The agent puts them together.
Its main routine is the loop from Chapter~1:

\begin{shell}
add the task to the conversation
repeat
    reply := the model's answer to the conversation
    add the reply to the conversation
    if the reply asks for tools then
        run them, and add their results to the conversation
    else
        the task is done
\end{shell}

Before the agent itself, three small things it needs: someone to decide
whether a tool may run, a tally of the tokens used, and a way to show what
is going on.

\section{Who decides}

Whether a tool that changes things may run is decided by a
\code{Permission}. Chapter~\ref{ch:permission} has the interesting ones,
including asking you. For now we need just the simplest, which allows
everything:

\begin{newfile}{permissions.nex}
-- permissions.nex: who decides whether a tool may run.
--
-- Reading the workspace is always allowed: it can't do any harm. A tool
-- that changes things (editing a file, running a command) runs only if
-- the policy allows it. When it doesn't, the model is told so, and can
-- try something else or ask.

intern Tools

deferred class Permission
feature
  -- Whether `tool` may run on `input`. It is only asked about calls
  -- whose input has been checked.
  allows(tool: Tool, input: Json_Value): Boolean
  require
    checked_input: tool.problem_with(input) = ""
  deferred
end

-- Everything is allowed: `--yes`. For a workspace you can throw away.
class Allow_Everything inherit Permission
create
  make() do
  end

feature
  allows(tool: Tool, input: Json_Value): Boolean do
    result := true
  end
end
\end{newfile}

\section{Counting tokens}

A real model's replies say how many tokens each request read and wrote
(Chapter~\ref{ch:conversation}). The agent adds them up, so that it can
tell you what a task used. \code{grouped} writes big numbers the readable
way, with commas.

\begin{newfile}{context.nex}
-- context.nex: keeping count of tokens, and starting afresh when the
-- conversation grows too long.
--
-- Every request sends the whole conversation, so each one is longer
-- than the last. The model can only read so much (its context window),
-- and every token read costs money and time. When the conversation
-- passes a budget, the agent asks the model to sum up what it has done,
-- and carries on in a new conversation that starts from that summary.

intern Messages
intern Helpers

-- What a session has used so far.
class Usage
create
  make() do
  end

feature
  replies: Integer        -- how many replies
  read: Integer           -- tokens read, over all requests
  cached: Integer         -- how many of those came from the cache
  written: Integer        -- tokens written

  add(reply: Reply) do
    replies := replies + 1
    read := read + reply.context_tokens()
    cached := cached + reply.cached_tokens()
    written := written + reply.output_tokens()
  end

  -- "12 replies; read 45,210 tokens (38,000 from the cache), wrote 1,234".
  to_string(): String do
    result := count_text(replies, "reply")
    if read > 0 then
      result := result + "; read " + grouped(read) + " tokens"
      if cached > 0 then
        result := result + " (" + grouped(cached) + " from the cache)"
      end
      result := result + ", wrote " + grouped(written)
    end
  end
end

-- A whole number with its thousands grouped: grouped(45210) is "45,210".
function grouped(n: Integer): String do
  let digits: String := n.abs.to_string
  result := ""
  let i: Integer := digits.length
  from until i <= 3 do
    result := "," + digits.substring(i - 3, i) + result
    i := i - 3
  end
  result := digits.substring(0, i) + result
  if n < 0 then
    result := "-" + result
  end
end
\end{newfile}

Note what \code{read} adds up. Every request sends the whole conversation,
so the tokens read grow with every reply: the tenth request reads
everything the first nine did, and more. That is why caching matters
(Chapter~\ref{ch:claude}), and why long conversations need care
(Chapter~\ref{ch:long}).

\section{How a task went}

When the agent stops working on a task, it says why. The model may have
finished (the usual case), or the agent may have stopped it: after too
many replies, because a reply was cut off, or because the model declined
the task. A scripted model can also run out of lines.

\begin{newfile}{agent.nex}
-- agent.nex: the agent loop.
--
-- Send the conversation to the model. If the reply asks for tools, run
-- them, add their results to the conversation, and send it again. When
-- a reply asks for nothing more, the task is done.

intern Model
intern Tools
intern Permissions
intern Context
intern Helpers

-- How a task went.
class Run_Result
create
  make(ending, text: String, replies, calls: Integer) do
    this.ending := ending
    this.text := text
    this.replies := replies
    this.calls := calls
  end

feature
  -- "done": the model finished; "turn limit": it was still going when
  -- the agent stopped it; "cut off": its last reply was too long;
  -- "refused": it declined; "out of script": a recorded or scripted
  -- model had no more replies.
  ending: String
  text: String      -- what the model said last
  replies: Integer  -- how many replies the task took
  calls: Integer    -- how many tool calls

  finished(): Boolean do
    result := ending = "done"
  end
end
\end{newfile}

\section{Showing what happens}

An agent that works in silence is hard to trust and impossible to debug.
The \code{Reporter} prints each step as it happens: the task, what the
model says, each tool call, and the first line of each result. It can also
be told to print the results in full, or to print nothing at all (for our
checks, in Chapter~\ref{ch:checks}).

\begin{append}{agent.nex}
-- Prints what the agent does, as it does it.
class Reporter
create
  -- `quiet` prints nothing; `full_output` prints every tool's output in
  -- full, not just its first line.
  make(quiet, full_output: Boolean) do
    this.quiet := quiet
    this.full_output := full_output
  end

feature
  quiet: Boolean
  full_output: Boolean

  on_task(task: String) do
    say("\n> " + task)
  end

  on_text(text: String) do
    say("\n" + text)
  end

  on_call(description: String) do
    say("  " + description)
  end

  on_result(text: String, is_error: Boolean) do
    if full_output then
      say(indented(text.trim(), "    | "))
    else
      let lines: Integer := lines_of(text).length
      say("    " + (when is_error then "error: " else "" end) + preview(text, 70)
        + (when lines > 1 then "  (" + count_text(lines, "line") + ")" else "" end))
    end
  end

  on_finish(outcome: Run_Result, usage: Usage) do
    let how: String := when outcome.finished() then "done" else "stopped: " + outcome.ending end
    say("\n(" + how + " after " + count_text(outcome.replies, "reply") + " and "
      + count_text(outcome.calls, "tool call") + ". So far: " + usage.to_string + ")")
  end

  private feature
    say(text: String) do
      if not quiet then
        console.print_line(text)
      end
    end
end
\end{append}

\section{The system prompt}

The \emph{system prompt} is the standing instructions, sent with every
request. Ours says what the model is doing and where, and how to go about
it: look before changing anything, make small changes, check them, and say
so when stuck. These are the habits of a careful programmer, and a model
follows them much better when they are written down. It ends by saying
what Nex is, since the model may never have seen it.

\begin{append}{agent.nex}
-- The standing instructions: who the model is in this conversation, and
-- how to work.
function system_prompt(): String do
  result := "You are a coding agent working in a project folder (the workspace). "
    + "Use the tools to read and search its files, change them, and run commands in it. "
    + "All paths are relative to the workspace.\n\n"
    + "Work the way a careful programmer would:\n"
    + "- Before changing anything, look around: read the files that matter, and run the "
    + "project's checks or tests to see where things stand.\n"
    + "- Make the smallest change that does the job, in the style of the code around it.\n"
    + "- After a change, run the checks again to see that it worked and broke nothing else.\n"
    + "- If something you need to know is unclear, or a problem is not one you can fix, "
    + "stop and say so instead of guessing.\n\n"
    + "When you have finished, reply without calling a tool: say in a few sentences what "
    + "was wrong, what you changed, and how you checked it.\n\n"
    + "The projects you work on may be written in Nex, a language in the Eiffel tradition: "
    + "classes, `require`/`ensure` contracts, `--` comments. Run a Nex program with "
    + "`nex file.nex`."
end

-- Whether no two tools share a name.
function names_differ(tools: Array[Tool]): Boolean do
  let names: Array[String] := []
  result := true
  across tools as tool do
    if names.contains(tool.name()) then
      result := false
    end
    names.add(tool.name())
  end
end
\end{append}

\section{The agent}

The agent is made from a model, its tools, a permission and a reporter. It
keeps the conversation (over every task you give it, so that a second task
can refer to the first) and a tally of tokens. Two tools with the same name
would make the model's requests ambiguous, so that is ruled out at the
start.

\begin{append}{agent.nex}
class Agent
create
  make(model: Model, tools: Array[Tool], permission: Permission, reporter: Reporter)
  require
    distinct_names: names_differ(tools)
  do
    this.model := model
    this.tools := tools
    this.permission := permission
    this.reporter := reporter
    conversation := create Conversation.make()
    usage := create Usage.make()
    system := system_prompt()
    max_replies := 40
  end

feature
  model: Model
  tools: Array[Tool]
  permission: Permission
  reporter: Reporter
  conversation: Conversation   -- everything so far, over every task
  usage: Usage
  system: String
  max_replies: Integer         -- the most replies one task may take

  set_max_replies(n: Integer)
  require
    positive: n > 0
  do
    max_replies := n
  end
\end{append}

\code{max_replies} is a safety net. A model that keeps calling tools
without getting anywhere (it happens: a test that can't be made to pass, a
misunderstanding that never clears up) would otherwise go on, at a cost,
for ever.

Now the loop. It is the plan above, with every way out made explicit:

\begin{continued}{agent.nex}
  -- Work on `task` until the model says it is done, or the agent stops
  -- it. The conversation carries on from earlier tasks, so a task can
  -- refer to one before it ("now add a check for that").
  run(task: String): Run_Result
  require
    something_to_do: task.trim() /= ""
  do
    reporter.on_task(task)
    conversation.add_user_text(task)
    let replies: Integer := 0
    let calls: Integer := 0
    let last_words: String := ""
    let ending: String := ""
    from until ending /= "" do
      if not model.can_reply() then
        ending := "out of script"
      elseif replies >= max_replies then
        ending := "turn limit"
      else
        let reply: Reply := model.reply(create Request.make(system, conversation, definitions()))
        replies := replies + 1
        usage.add(reply)
        conversation.add_reply(reply)
        if reply.text() /= "" then
          reporter.on_text(reply.text())
          last_words := reply.text()
        end
        let wanted: Array[Tool_Call] := reply.tool_calls()
        if not wanted.is_empty then
          let results: Array[Tool_Result] := []
          across wanted as call do
            results.add(execute(call, reply.stop_reason() = "max_tokens"))
            calls := calls + 1
          end
          conversation.add_tool_results(results)
        else
          ending := ending_for(reply.stop_reason())
        end
      end
    end
    result := create Run_Result.make(ending, last_words, replies, calls)
    reporter.on_finish(result, usage)
  ensure
    no_calls_left_unanswered: conversation.waiting.is_empty
  end
\end{continued}

A few things to notice.

\begin{itemize}
\item The loop decides what to do from the reply's content, not just its
  \code{stop_reason}: if there are tool calls, they are answered. That
  matters in one case. A reply can be cut off part way, when it reaches the
  longest reply the model is allowed (\code{stop_reason} is
  \code{"max_tokens"}). If it was in the middle of a tool call, that call's
  input may be incomplete, so \code{execute} is told, and refuses to run
  it. But the call is still answered, with an error, because every call
  must be answered before the conversation can go on.
\item All the results of one reply go back in one message. A model can ask
  for several tools at once (to read three files, say), and answering them
  together is what the API expects.
\item The postcondition says that however the task ended, no tool call is
  left unanswered. That is what lets you give the agent another task
  afterwards.
\end{itemize}

\section{Running one tool call}

\code{execute} is where the agent's rules are applied. In order: the tool
must exist, the reply must not have been cut off, the input must be valid,
and the permission must allow it. Only then does the tool run, and even
then, if it fails with an exception, the failure becomes an error result.
\emph{Nothing a model asks for can stop the agent.} Whatever goes wrong
becomes a message the model can read, and act on.

\begin{continued}{agent.nex}
  -- Run one tool call and say how it went. Whatever goes wrong becomes
  -- a result the model can read; the agent itself carries on.
  execute(call: Tool_Call, cut_off: Boolean): Tool_Result do
    let outcome: Tool_Outcome := create Tool_Outcome.error("")
    let tool: ?Tool := tool_named(call.name)
    if ?tool as t then
      reporter.on_call(t.describe(call.input))
      let problem: String := t.problem_with(call.input)
      if cut_off then
        outcome := create Tool_Outcome.error("Your reply was cut off before this call was "
          + "complete, so it was not run. Try again, sending less in one call.")
      elseif problem /= "" then
        outcome := create Tool_Outcome.error("Not run: " + problem + ".")
      elseif not permission.allows(t, call.input) then
        outcome := create Tool_Outcome.error("The user did not allow this, so it was not run. "
          + "Don't try to get the same effect another way; if you need it, say why.")
      else
        do
          outcome := t.run(call.input)
        rescue
          outcome := create Tool_Outcome.error("The tool failed: " + exception)
        end
      end
    else
      reporter.on_call(call.name + " ?")
      outcome := create Tool_Outcome.error("There is no tool called `" + call.name
        + "`. The tools are: " + ", ".join(tool_names()) + ".")
    end
    reporter.on_result(outcome.text, outcome.is_error)
    result := create Tool_Result.make(call.id, outcome.text, outcome.is_error)
  ensure
    answers_the_call: result.call_id = call.id
  end
\end{continued}

The message for a refusal is worth a second look: ``Don't try to get the
same effect another way.'' Without it, a model refused an edit will
sometimes try writing the whole file instead, or running \code{sed}: it is
trying to be helpful, and the refusal looks to it like an obstacle. Saying
plainly what a refusal means is part of the job.

The \code{do} \ldots\ \code{rescue} \ldots\ \code{end} block is how Nex
catches an exception: if anything in the \code{do} part fails, the
\code{rescue} part runs instead of the program stopping.

The rest is bookkeeping: the tools' definitions for the request, and
finding a tool by name.

\begin{continued}{agent.nex}
  -- The tools as the API wants them described.
  definitions(): Array[Map[String, Any]] do
    across tools as tool do
      result.add(tool.definition())
    end
  end

  tool_named(name: String): ?Tool do
    across tools as tool do
      if tool.name() = name then
        result := tool
      end
    end
  end

  tool_names(): Array[String] do
    across tools as tool do
      result.add(tool.name())
    end
  end

  private feature
    -- What a reply that asks for no tools means for the task.
    ending_for(stop_reason: String): String do
      if stop_reason = "max_tokens" then
        result := "cut off"
      elseif stop_reason = "refusal" then
        result := "refused"
      else
        result := "done"     -- "end_turn", or a stop sequence
      end
    end
end
\end{continued}

\code{tool_named} returns a \code{?Tool}, which is \code{nil} when there is
no such tool; \code{execute} checks with \code{if ?tool as t}.

\section{A first run}

Our agent can't do much yet: it can only look. But it is complete, so
let's give it a script. The script lists the files, asks for a tool that
doesn't exist, gets a parameter wrong, reads a file, and finishes:

\begin{tryit}{try_agent.nex}
intern Agent

let ws: Workspace := fresh_workspace(user_path("sample_project"), user_path("scratch"))
let tools: Array[Tool] := [create List_Files.make(ws), create Read_File.make(ws)]
let model: Scripted_Model := create Scripted_Model.make([
  use_tool("Let me see what the project has.", "list_files", '{}'),
  use_tool("", "delete_everything", '{}'),
  use_tool("", "read_file", '{"name": "fines.nex"}'),
  use_tool("", "read_file", '{"path": "fines.nex"}'),
  say("A Fine_Policy charges so much per day late, up to a maximum.")
])
let agent: Agent := create Agent.make(model, tools, create Allow_Everything.make(),
  create Reporter.make(false, false))
let outcome: Run_Result := agent.run("What does fines.nex do?")
console.print_line("ending: " + outcome.ending)
console.print_line("what the model was told about delete_everything: " + model.sent.get(2))
\end{tryit}
\tryoutput{try_agent}

The mistakes didn't stop anything: each became an error the model could
read, as the last line shows. A script can't learn from an error, of
course, but a real model does, as we will see.
